\section{Conclusion}
\label{sec:conclusion}

When a false step sits in a model's own reasoning, what happens next depends on what checking would cost. A value the model can reread is caught in proportion to the model's strength, from less than half the time for the small open-weight models to almost always at the frontier. A value the model would have to recompute is kept. It is absorbed into the continuation at 88 to 100 percent by every model with enough solved programs to measure, unmoved by scale, by showing the operands beside the wrong result, or by four escalating instructions to check. For Haiku 4.5 it is switched on by a single operation of verification depth. Five years of model progress built near-perfect rereaders and did not produce a recomputer. Until training does, systems that chain model reasoning should treat every committed computed value as unverified input, because the model treats it as settled truth.
